\section{Medição de Parâmetros}


\subsection{Resistência}
A resistência de linha foi obtida pelo método de quatro terminais (Kelvin),
aplicando corrente entre dois terminais e medindo a queda de tensão sobre $A$ e
$B$, com $R_{eq} = V_{AB}/I$.
Obteve-se $R_{LL} = 17{,}37\ \text{m}\Omega$
(desvio padrão $0{,}02\ \text{m}\Omega$), correspondendo a
$R_{Y} = R_{LL}/2 = 8{,}69\ \text{m}\Omega$ no modelo em estrela equivalente.

\begin{table}[htb]
    \centering
    \caption{Medição da resistência equivalente de linha do motor.}
    \label{tab:resistencia_motor}
    \sisetup{output-decimal-marker={,}}
    \begin{tabular}{
        S[table-format=2.3]
        S[table-format=1.3]
        S[table-format=2.2]
    }
        \toprule
        {$V_{AB}$ [\si{\milli\volt}]} & {$I$ [\si{\ampere}]} & {$R_{eq}$ [\si{\milli\ohm}]} \\
        \midrule
        3,654  & 0,210 & 17,40 \\
        7,172  & 0,412 & 17,41 \\
        9,629  & 0,554 & 17,38 \\
        11,407 & 0,657 & 17,36 \\
        12,714 & 0,732 & 17,37 \\
        14,314 & 0,824 & 17,37 \\
        16,145 & 0,930 & 17,36 \\
        18,150 & 1,046 & 17,35 \\
        19,871 & 1,145 & 17,35 \\
        20,810 & 1,199 & 17,36 \\
        22,645 & 1,306 & 17,34 \\
        24,715 & 1,426 & 17,33 \\
        \midrule
        \multicolumn{2}{r}{Média $R_{eq}$} & 17,37 \\
        \multicolumn{2}{r}{Desvio padrão} & 0,02 \\
        \multicolumn{2}{r}{$R_{Y}=R_{LL}/2$} & 8,69 \\
        \bottomrule
    \end{tabular}
\end{table}

\subsection{Indutância}
A indutância foi determinada pelo método da ressonância, associando capacitores
de capacitância conhecida em paralelo aos terminais do motor e excitando o
circuito $LC$ com um gerador em modo de varredura (\textit{AC sweep}).

Na ressonância, $f_{res} = 1/(2\pi\sqrt{LC})$, de onde $L = 1/[(2\pi f_{res})^2 C]$.
A frequência de cada pico foi refinada manualmente. O uso de diferentes
capacitores, medidos individualmente, fornece estimativas independentes de $L$
para o mesmo conjunto, e a convergência entre elas valida a medição. Obteve-se $L_{LL} = 13{,}90 \pm 0{,}14\ \mu\text{H}$ e
$L_{Y} = L_{LL}/2 = 6{,}95 \pm 0{,}07\ \mu\text{H}$.

\begin{table}[H]
    \centering
    \caption{Determinação da indutância do motor pelo método da ressonância,
    para diferentes valores de capacitância.}
    \label{tab:indutancia_motor}
    \sisetup{output-decimal-marker={,}}
    \begin{tabular}{
        S[table-format=1.3]
        S[table-format=2.1]
        S[table-format=2.2]
        S[table-format=1.2]
    }
        \toprule
        {$C$ [\si{\micro\farad}]} & {$f_{res}$ [\si{\kilo\hertz}]} &
        {$L_{LL}$ [\si{\micro\henry}]} & {$L_{Y}$ [\si{\micro\henry}]} \\
        \midrule
        2,200 & 28,9 & 13,79 & 6,89 \\
        4,443 & 20,3 & 13,84 & 6,92 \\
        6,700 & 16,4 & 14,06 & 7,03 \\
        \midrule
        \multicolumn{2}{r}{Média} & {\SI{13,90}{\micro\henry}} & {\SI{6,95}{\micro\henry}} \\
        \multicolumn{2}{r}{Desvio padrão} & {\SI{0,14}{\micro\henry}} & {\SI{0,07}{\micro\henry}} \\
        \bottomrule
    \end{tabular}
\end{table}


